\section{Compact completion of the full averaged residual}
\label{mcstress:section}

This section computes a compact symmetric stress and a compact pressure
for the complete averaged residual of the actual finite curl candidate.
It retains the original radial bumps depending on $q(z,t)$, both axial
fluxes, the radial force, and the full cylindrical connection terms.
The result is a stress--pressure representation of a computed force;
it does not by itself modify a velocity or identify the constructed
stress with a covariance of additional waves.

\subsection{The actual extended fields and all three averaged components}

Use right-handed cylindrical coordinates
$x=(r\cos\theta,r\sin\theta,z)$, with oriented orthonormal basis
$(e_r,e_\theta,e_z)$ and $r>0$. The retained finite source state is
$(u_*,p_*)$; the actual added field and its selected pressure are
$(v,\pi)=(w_y,\pi_y)$. On a finite collection of source labels the
added potential has the specified form
\begin{equation}\label{mcstress:actual_potential}
\begin{gathered}
 \mathcal A_y=\sum_j\operatorname{Re}
   \left[Q_j^{1/2-A}
     \frac{i\,n_j\times t_j}{k_jm_j|n_j|^2}
                      e^{ik_jm_j\Phi_j}\right],\qquad
 v=\operatorname{curl}_{\mathscr D}\mathcal A_y,\\
 t_j=\chi_jB_jy_3(t_{c,j}),\qquad
 y_3=\binom{\Theta_3}{\Omega_3},\qquad
 t_{c,j}=t_b+\alpha_jv_j,\qquad
 n_j^TB_j=0,\qquad k_jm_jp_j\in\mathbb Z\setminus\{0\}.
\end{gathered}
\end{equation}
Here $y_3$ is the actual signed third-growth trajectory, the cutoffs
and source frame are the retained functions, and $Q_j$ is constant
on each source chart. The chosen pressure is the real harmonic sum
with chart coefficient
\[
 Q_j^{-2A}
 \frac{i\{n_j^T\mathsf K_jt_j-(n'_j)^Tt_j\}}
      {k_jm_j|n_j|^2}e^{ik_jm_j\Phi_j}.
\]
All coefficient components are independent of $\theta$; their angular
dependence is the displayed nonzero integer harmonic. Their radial
and axial supports lie in a fixed compact set with $r$ bounded away
from zero on the finite time interval under consideration. Smooth
zero extension at each source cutoff is retained.

The source absolute auxiliary coordinate is $Y\in\mathbb T^2$ and
the physical evaluation is
\begin{equation}\label{mcstress:extended_derivatives}
\begin{gathered}
 Y(r,t)=v_r r^{d_r}+v_t t\pmod{\mathbb Z^2},\qquad
 v_r=(1,-b_g),\quad v_t=(b_g,1),\quad b_g=\sqrt2-1,\\
 \rho_g=\frac{\log(4-\sqrt2)}{\log(4+\sqrt2)},\quad
 \kappa_s=10^{-5},\quad d_r=2((1+h)\rho_g-h\kappa_s),\\
 \mathscr D_r=\partial_r+d_r r^{d_r-1}v_r\cdot\partial_Y,
 \qquad \mathscr D_t=\partial_t+v_t\cdot\partial_Y,
 \qquad \mathscr D_z=\partial_z.
\end{gathered}
\end{equation}
The angular derivative is $\partial_\theta$.
The exponent $h$ here is the unchanged source chart exponent. These
lifted derivatives commute, obey the product rule, and satisfy
$\mathscr D_r r=1$. Equivalently in Cartesian coordinates they are
$\partial_{x_i}+\sum_\alpha(\partial_{x_i}Y_\alpha)\partial_{Y_\alpha}$;
their commutators vanish because mixed derivatives of $Y$ commute.
Thus evaluation at the physical map preserves every differentiated
identity below.

The actual finite source has its retained representation
$u_* =\operatorname{curl}_{\mathscr D}\mathcal A_*+B_*e_\theta$,
with $\partial_\theta B_*=0$. Commutation of the lifted derivatives
proves $\operatorname{div}_{\mathscr D}\operatorname{curl}_{\mathscr D}=0$.
Also $\operatorname{div}_{\mathscr D}(B_*e_\theta)=0$.
Consequently both $u_*$ and $v$ are divergence free at every independent
$Y$, not only after physical evaluation.

Let the normalized averages act on cylindrical components before
evaluation:
\[
 \langle f\rangle_\theta=\frac1{2\pi}\int_0^{2\pi}f\,d\theta,
 \quad \langle f\rangle_Y=\int_{\mathbb T^2}f\,dY,
 \quad \langle f\rangle=\langle\langle f\rangle_\theta\rangle_Y,
 \quad\int_{\mathbb T^2}1\,dY=1.
\]
Every cylindrical curl operation retains the nonzero angular harmonic
in \eqref{mcstress:actual_potential}. For example its radial and axial
components are $r^{-1}\partial_\theta\mathcal A_z-\partial_z\mathcal A_\theta$
and $(\mathscr D_r+r^{-1})\mathcal A_\theta-r^{-1}\partial_\theta\mathcal A_r$.
Integrating the integer harmonic proves
\begin{equation}\label{mcstress:zero_mean}
 \langle v\rangle_\theta=0,\qquad
 \langle\pi\rangle_\theta=0
 \quad\hbox{at every }(r,z,t,Y).
\end{equation}
This statement includes the complete coefficient curl and every cutoff
derivative, rather than only the principal amplitude.

For a smooth periodic coefficient $f$, integration of its auxiliary
derivatives gives
\begin{equation}\label{mcstress:average_derivatives}
 \langle\mathscr D_r f\rangle=\partial_r\langle f\rangle,
 \quad \langle\mathscr D_t f\rangle=\partial_t\langle f\rangle,
 \quad \langle\partial_z f\rangle=\partial_z\langle f\rangle,
 \quad \langle\partial_\theta f\rangle=0.
\end{equation}
Applying the first identity twice also treats $\mathscr D_r^2$;
the derivative of its variable coefficient has zero auxiliary average
because that coefficient is independent of $Y$.

At arbitrary fixed physical viscosity $\nu_{\rm NS}>0$, the exact
residual increment is
\begin{equation}\label{mcstress:residual_increment}
 \Delta R=\mathscr D_t v+(u_*\cdot\nabla_{\mathscr D})v
 +(v\cdot\nabla_{\mathscr D})u_*
 +(v\cdot\nabla_{\mathscr D})v
 -\nu_{\rm NS}\Delta_{\mathscr D}v+\nabla_{\mathscr D}\pi.
\end{equation}
This is the residual of the actual modified state minus the source
state's residual at the same viscosity. Put
\begin{equation}\label{mcstress:actual_tensor}
 \bar u_* =\langle u_*\rangle_\theta,\qquad
 w_*=u_*-\bar u_*,\qquad
 S=\langle w_*\otimes v+v\otimes w_*+v\otimes v\rangle.
\end{equation}
It is a smooth real symmetric axisymmetric tensor compactly supported
in the annular and axial set at issue. Each summand contains a factor
of the actual compact field $v$.

\begin{theorem}\label{mcstress:full_mean}
The complete averaged residual is the full Cartesian divergence of
the actual tensor $S$:
\begin{equation}\label{mcstress:Ffull}
 \langle\Delta R\rangle
 =F_r e_r+F_2e_\theta+F_1e_z
 =\operatorname{div}S,
\end{equation}
where
\begin{equation}\label{mcstress:all_components}
\begin{aligned}
 F_r&=(\partial_r+r^{-1})S_{rr}-r^{-1}S_{\theta\theta}
                                      +\partial_zS_{rz},\\
 F_2&=(\partial_r+2r^{-1})S_{r\theta}+\partial_zS_{z\theta},\\
 F_1&=(\partial_r+r^{-1})S_{rz}+\partial_zS_{zz}.
\end{aligned}
\end{equation}
\end{theorem}
\begin{proof}
For divergence-free fields, direct Cartesian component expansion
rewrites the three advection terms of \eqref{mcstress:residual_increment}
as the lifted divergence of
$u_*\otimes v+v\otimes u_*+v\otimes v$.
The mean-flow cross tensor has zero angular average by
\eqref{mcstress:zero_mean}, even when $\bar u_*$ depends on $Y$.
Thus its averaged divergence is zero by
\eqref{mcstress:average_derivatives} and the cylindrical divergence
formula.

To check the linear radial component explicitly, let
\[
 \mathscr L_0=\mathscr D_r^2+r^{-1}\mathscr D_r
             +r^{-2}\partial_\theta^2+\partial_z^2.
\]
The complete cylindrical vector Laplacian is
\[
 \begin{aligned}
 (\Delta_{\mathscr D}v)_r&=(\mathscr L_0-r^{-2})v_r
                       -2r^{-2}\partial_\theta v_\theta,\\
 (\Delta_{\mathscr D}v)_\theta&=(\mathscr L_0-r^{-2})v_\theta
                       +2r^{-2}\partial_\theta v_r,\qquad
 (\Delta_{\mathscr D}v)_z=\mathscr L_0v_z.
 \end{aligned}
\]
Each average is zero by \eqref{mcstress:zero_mean} and
\eqref{mcstress:average_derivatives}. The radial pressure average is
$\langle\mathscr D_r\pi\rangle=\partial_r\langle\pi\rangle=0$;
the angular and axial components vanish as well. The averaged time
derivative is zero. These verifications retain the radial connection
term and physical viscosity before taking the average.

Finally, for a symmetric tensor independent of $\theta$ in cylindrical
components, differentiation of
$e_r=(\cos\theta,\sin\theta,0)$ and
$e_\theta=(-\sin\theta,\cos\theta,0)$ gives exactly the three
expressions in \eqref{mcstress:all_components}. One direct verification
is to write its Cartesian matrix as $O_\theta S O_\theta^T$, where
$O_\theta=[e_r,e_\theta,e_z]$, and use
$\partial_x=\cos\theta\partial_r-r^{-1}\sin\theta\partial_\theta$
and $\partial_y=\sin\theta\partial_r+r^{-1}\cos\theta\partial_\theta$.
The derivatives of $O_\theta$ give $-S_{\theta\theta}/r$ in the radial
row, $2S_{r\theta}/r$ in the angular row, and $S_{rz}/r$ in the axial
row. Averaging the lifted version replaces $\mathscr D_r$ by
$\partial_r$. This proves the full Cartesian identity.
\end{proof}

\subsection{The compact fluxes and the original varying radial bumps}

For $e=1,2$, write $\tau_1=z$, $\tau_2=\theta$, and set
\begin{equation}\label{mcstress:fluxes}
 M_e(z,t)=\int_0^\infty r^eF_e(r,z,t)\,dr,
 \qquad
 J_e(z,t)=\int_0^\infty r^eS_{z\tau_e}(r,z,t)\,dr.
\end{equation}
Multiplication of the corresponding row of
\eqref{mcstress:all_components} by $r^e$ makes its radial term
$\partial_r(r^e S_{r\tau_e})$. The radial boundary values vanish
by compact annular support. Therefore
\begin{equation}\label{mcstress:flux_identity}
 M_e=\partial_zJ_e,\qquad
 \int_{\mathbb R}M_e(z,t)\,dz=0\quad(e=1,2).
\end{equation}
Both $J_e$ are smooth and compactly supported in $z$. No claim that
either $M_e$ vanishes pointwise is used.

Retain the original $q$-dependent bumps, for $q(z,t)>0$ on the compact
source patch:
\begin{equation}\label{mcstress:bumps}
 b_e(r,z,t)=q(z,t)^{-(e+1)/2}
             \widehat b_e(r/\sqrt{q(z,t)}),\qquad
 \widehat b_e\in C_c^\infty((0,\infty)),\qquad
 \int_0^\infty x^e\widehat b_e(x)\,dx=1.
\end{equation}
The fixed profiles may be taken to be the already selected ones.
The change of variable $r=\sqrt q\,x$ proves
$\int r^e b_e\,dr=1$ for every $(z,t)$. Choose a fixed annulus
containing the supports of $S$ and of the bumps on the relevant compact
$(z,t)$ set. This is possible because $q$ has positive lower and finite
upper bounds there. Outside the compact axial set where a flux or
force is nonzero all the formulas below extend by zero.

The exact derivatives and cumulative bump are
\begin{equation}\label{mcstress:bump_derivatives}
\begin{gathered}
 \partial_zb_e=-\frac{q_z}{2q}
                  \bigl((e+1)b_e+r\partial_rb_e\bigr),\\
 \beta_e(r,z,t)=\int_0^r s^eb_e(s,z,t)\,ds
              =\int_0^{r/\sqrt q}x^e\widehat b_e(x)\,dx,\\
 \partial_z\beta_e=-\frac{q_z}{2q}r^{e+1}b_e(r,z,t),\qquad
 \int_0^\infty r^e\partial_zb_e\,dr=0.
\end{gathered}
\end{equation}
The first identity is the product and chain rule. The second is the
same radial substitution as the normalization. Differentiating its
upper limit proves the third. Differentiating the normalization under
the uniformly compact radial integral proves the fourth. In particular
none of these identities assigns a sign to $q_z$.

\subsection{Corrected radial primitives and the complete symmetric tensor}

Define the actual corrected primitives
\begin{equation}\label{mcstress:sigma}
 \sigma_e(r,z,t)=-r^{-e}\int_0^r s^e
       \bigl(F_e(s,z,t)-\partial_z(b_e(s,z,t)J_e(z,t))\bigr)\,ds,
 \qquad e=1,2.
\end{equation}
Their exact radial equations are
\begin{equation}\label{mcstress:sigma_equation}
 (\partial_r+e/r)\sigma_e
      =-F_e+\partial_z(b_eJ_e).
\end{equation}
This follows by multiplying \eqref{mcstress:sigma} by $r^e$ and
using the fundamental theorem of calculus.

For comparison with the initially proposed primitive, put
\[
 \sigma_e^{\rm old}=-r^{-e}\int_0^r s^e
                       (F_e-b_e\partial_zJ_e)\,ds.
\]
Every missing term can be displayed exactly:
\begin{equation}\label{mcstress:primitive_correction}
 \sigma_e=\sigma_e^{\rm old}+r^{-e}J_e\partial_z\beta_e
         =\sigma_e^{\rm old}-\frac{rq_z}{2q}b_eJ_e.
\end{equation}
Indeed expanding $\partial_z(b_eJ_e)$ in
\eqref{mcstress:sigma} proves the first equality, and
\eqref{mcstress:bump_derivatives} proves the second. With the old
primitive the angular or axial stress divergence below would instead
be $-F_e-(\partial_zb_e)J_e$. Thus the correction is necessary for the
retained varying bump. It vanishes in the special case $q_z=0$.

There is an exact expression entirely in the original tensor entries:
\begin{equation}\label{mcstress:compact_flux_correction}
\begin{gathered}
 C_e(r,z,t)=\int_0^r s^eS_{z\tau_e}(s,z,t)\,ds
                                    -\beta_e(r,z,t)J_e(z,t),\\
 \sigma_e=-S_{r\tau_e}-r^{-e}\partial_zC_e.
\end{gathered}
\end{equation}
To prove it, integrate the corresponding component of
\eqref{mcstress:all_components} only up to $r$, obtaining
$\int_0^r s^eF_e\,ds=r^eS_{r\tau_e}
+\partial_z\int_0^r s^eS_{z\tau_e}\,ds$. Substitute this identity
and $\int_0^r s^e b_e\,ds=\beta_e$ into
\eqref{mcstress:sigma}. Below both radial supports $C_e=0$;
above both supports its two summands are $J_e$ and $J_e$, so it is
again zero. This proves its compact support as well as the exact
relation of the constructed radial stress to the actual covariance
tensor.

\begin{theorem}\label{mcstress:completion}
The real symmetric axisymmetric tensor with cylindrical entries
\begin{equation}\label{mcstress:tensor}
\begin{gathered}
 T_{rr}=0,\qquad T_{r\theta}=T_{\theta r}=\sigma_2,\qquad
 T_{rz}=T_{zr}=\sigma_1,\\
 T_{z\theta}=T_{\theta z}=-b_2J_2,\qquad
 T_{zz}=-b_1J_1,\qquad
 T_{\theta\theta}=rF_r+r\partial_z\sigma_1
\end{gathered}
\end{equation}
is smooth and compactly supported in an annular cylinder and obeys
the full Cartesian identity
\begin{equation}\label{mcstress:divergence}
 \operatorname{div}T=-(F_r e_r+F_2e_\theta+F_1e_z).
\end{equation}
If the term $rF_r$ is omitted only in $T_{\theta\theta}$, the exact
divergence is instead $-(0,F_2,F_1)$ in cylindrical components.
\end{theorem}
\begin{proof}
The integral of the numerator in \eqref{mcstress:sigma} over the
entire radial half-line is
\[
 M_e-\partial_z\left(J_e\int_0^\infty s^eb_e\,ds\right)
  =\partial_zJ_e-\partial_zJ_e=0.
\]
Thus the numerator primitive is zero below its support and zero above
it. Division by $r^e$ introduces no singularity because it is zero on
a fixed neighborhood of the axis. All derivatives under its fixed
compact radial integral exist, and the normalization identity holds
smoothly in $(z,t)$, so its extension by zero is smooth. Compact axial
support follows because both $F_e$ and $J_e$ vanish outside the fixed
compact axial set. The same statements hold for all entries of
\eqref{mcstress:tensor}. Passing to the Cartesian matrix
$O_\theta T O_\theta^T$ preserves smoothness away from the axis; near
the axis it is identically zero. Hence this is a smooth compactly
supported Cartesian tensor.

Apply the full symmetric axisymmetric divergence formula already
proved in \eqref{mcstress:all_components}. Its three rows are
\[
 \begin{aligned}
 (\operatorname{div}T)_r
   &=-\frac{rF_r+r\partial_z\sigma_1}{r}
                                     +\partial_z\sigma_1=-F_r,\\
 (\operatorname{div}T)_\theta
   &=(\partial_r+2/r)\sigma_2-\partial_z(b_2J_2)=-F_2,\\
 (\operatorname{div}T)_z
   &=(\partial_r+1/r)\sigma_1-\partial_z(b_1J_1)=-F_1.
 \end{aligned}
\]
The last two identities are \eqref{mcstress:sigma_equation}, including
every derivative of $b_e$. This proves the Cartesian vector identity
and, on omitting $rF_r$, its stated radial-zero variant.
\end{proof}

\subsection{Compact pressure, tracefree stress, and the exact gauge relation}

The trace and its prescribed pressure are explicitly
\begin{equation}\label{mcstress:pressure}
 \operatorname{tr}T=rF_r+r\partial_z\sigma_1-b_1J_1,
 \qquad
 p_T=\frac13\bigl(rF_r+r\partial_z\sigma_1-b_1J_1\bigr),
 \qquad T^0=T-p_TI_3.
\end{equation}
Then $\operatorname{tr}T^0=0$, and every field is smooth and compactly
supported in the same annular cylinder. The exact full cancellation is
\begin{equation}\label{mcstress:tracefree_cancellation}
 F+\operatorname{div}T^0+\nabla p_T=0.
\end{equation}
To check its sign, for any scalar $p$ Cartesian differentiation gives
$\operatorname{div}(pI_3)=\nabla p$. Therefore
$\operatorname{div}T^0=\operatorname{div}T-\nabla p_T
=-F-\nabla p_T$, which proves the displayed equation.
In components the new diagonal entries are
\[
 T^0_{rr}=-p_T,\quad
 T^0_{\theta\theta}=rF_r+r\partial_z\sigma_1-p_T,\quad
 T^0_{zz}=-b_1J_1-p_T.
\]
The axial stress divergence alone changes by $-\partial_zp_T$;
the pressure gradient contributes $+\partial_zp_T$. Thus the combined
stress--pressure correction preserves the exact tangential equations.
There is no claim that their stress term alone stays unchanged.

More generally $T-pI_3$ with any smooth compact scalar $p$ satisfies
$F+\operatorname{div}(T-pI_3)+\nabla p=0$.
The choice in \eqref{mcstress:pressure} is the unique one making
this stress tracefree. The map from symmetric tensors to
tracefree-tensor--scalar pairs,
\[
 T\longmapsto\bigl(T-\tfrac13\operatorname{tr}T\,I_3,
                         \tfrac13\operatorname{tr}T\bigr),
\]
has exact inverse $(U,p)\mapsto U+pI_3$ when $\operatorname{tr}U=0$.
It intertwines $\operatorname{div}$ with
$(U,p)\mapsto\operatorname{div}U+\nabla p$. This supplies the full
map between the two presentations, including compact support and signs.

\subsection{All total force and torque components and the divergence image}

The Cartesian components of the averaged force are
\[
 F=(F_r\cos\theta-F_2\sin\theta,
       F_r\sin\theta+F_2\cos\theta,F_1).
\]
Integration in $\theta$ with the physical volume element
$dx=r\,dr\,d\theta\,dz$ gives
\begin{equation}\label{mcstress:force_torque}
\begin{gathered}
 \int_{\mathbb R^3}F\,dx
  =2\pi e_z\int_{\mathbb R}M_1(z,t)\,dz=0,\\
 \int_{\mathbb R^3}x\times F\,dx
  =2\pi e_z\int_{\mathbb R}M_2(z,t)\,dz=0.
\end{gathered}
\end{equation}
For the torque calculation, direct multiplication in the oriented
basis gives
\[
 x\times F=-zF_2e_r+(zF_r-rF_1)e_\theta+rF_2e_z.
\]
The first two terms integrate to zero in $\theta$, and the third
gives the second equation of \eqref{mcstress:force_torque}.
Thus all three force and all three torque components vanish. The last
equalities use the already proved compact flux identities
\eqref{mcstress:flux_identity}. The same torque vanishes about every
fixed origin $x_0$, because subtracting $x_0$ changes it by
$-x_0\times\int F\,dx=0$. On a finite axial slab $a<z<b$, the
axial force and torque are respectively
$2\pi(J_1(b,t)-J_1(a,t))$ and
$2\pi(J_2(b,t)-J_2(a,t))$. The quantities need not vanish on an
individual slice.

There is also a direct Cartesian stress verification. Since
$F_j=-\partial_kT_{jk}$ and $T$ is compact, integration by parts gives
\begin{equation}\label{mcstress:first_moments}
 \int F_j\,dx=0,\qquad
 \int x_iF_j\,dx=\int T_{ji}\,dx=-\int S_{ji}\,dx.
\end{equation}
Symmetry makes the antisymmetric part of the first-moment matrix zero,
which is precisely all three torque components. Its symmetric part
is not asserted to vanish. The second equality with $S$ uses
$F=\operatorname{div}S$. Equivalently
$\operatorname{div}(T+S)=0$; multiplying this equation by $x_i$ and
integrating proves $\int(T+S)_{ji}\,dx=0$. In the tracefree-pressure
presentation the same first moment is
$\int T^0_{ji}\,dx+\delta_{ij}\int p_T\,dx$.
All these identities persist under any finite time differentiation on
a common compact support, by differentiation under the integral.

For completeness, this construction proves the exact image of compact
symmetric-stress divergence in the axisymmetric class. Let
an open axial interval $I_z$ contain the actual axial supports, with
$q>0$ smooth on a neighborhood of its relevant compact subsets at the
fixed time. In the following two spaces all supports are compact
subsets of $\{r>0,\ z\in I_z\}$. Thus the original bump is defined
everywhere it is used, and integrating between axial support endpoints
stays inside $I_z$.
Let
$\mathcal F_0$ be the space of smooth real axisymmetric compact forces
supported away from $r=0$ and satisfying exactly
\[
 \int_{\mathbb R}\int_0^\infty rF_1\,dr\,dz=0,
 \qquad
 \int_{\mathbb R}\int_0^\infty r^2F_2\,dr\,dz=0.
\]
Let $\mathcal S_c$ be the space of smooth real compact symmetric
axisymmetric Cartesian tensors supported away from $r=0$. Define
$\mathcal D:\mathcal S_c\to\mathcal F_0$ by
$\mathcal D(U)=-\operatorname{div}U$.
Integration by parts proves that its values satisfy both defining
identities of $\mathcal F_0$, and axisymmetry makes the other force
and torque components vanish as above.

For any $F\in\mathcal F_0$, compute its actual radial moments and set
\[
 J_e(z)=\int_{-\infty}^z\int_0^\infty r^eF_e(r,s)\,dr\,ds.
\]
This is a smooth compact axial function: its derivative is $M_e$, it
vanishes below the axial support, and the defining zero total integral
makes it vanish above that support. Use the fixed bumps
\eqref{mcstress:bumps}, the corrected primitive
\eqref{mcstress:sigma}, and all entries of
\eqref{mcstress:tensor}. The proof of Theorem~\ref{mcstress:completion}
uses only these identities and proves a compact tensor
$\mathcal C_bF$ with
\begin{equation}\label{mcstress:right_inverse}
 \mathcal D\mathcal C_bF=F\qquad(F\in\mathcal F_0).
\end{equation}
Thus the two total moments are necessary and sufficient, with an
explicit right inverse. For the actual force in
\eqref{mcstress:Ffull}, this primitive in $z$ equals the already
computed $J_e$ from $S$, because both have derivative $M_e$ and vanish
at the negative axial end. No additional compatibility condition has
been imposed on the actual candidate.

For fixed $q$ and bump profiles, all operations defining
$\mathcal C_b$ are linear. Hence
$\mathcal P_b=\mathcal C_b\mathcal D$ is a projection on
$\mathcal S_c$: \eqref{mcstress:right_inverse} gives
$\mathcal P_b^2=\mathcal P_b$ and
$\ker\mathcal P_b=\ker\mathcal D$.
In particular the completed actual tensor is
$T=\mathcal P_b(-S)$, and $T+S$ is a compact divergence-free symmetric
tensor. This gives an exact relation to the actual averaged covariance,
not merely equality of its two radial moments.

\subsection{Exact source chart, bump, stress, and pressure powers}

Fix one source chart and retain all of its constants:
\begin{equation}\label{mcstress:chart}
 r=\sqrt Q\,R,\qquad z=Q^D Z,\qquad1-t=QT,
 \qquad A=\tfrac12+h,\quad D=\tfrac12-h,\quad
 \varepsilon=Q^h.
\end{equation}
The dyadic label $Q$ is fixed under differentiation. The averaged quantities
are independent of $Y$ and $\theta$, so in this
subsection the complete normalized spatial derivatives are
$\partial_R$, $R^{-1}\partial_\theta$ and $\varepsilon\partial_Z$,
with the same cylindrical basis derivatives as before.
Define the exact chart representatives
\begin{equation}\label{mcstress:chart_representatives}
\begin{gathered}
 S^*=Q^{2A}S,\qquad F^*=Q^{2A+1/2}F,\qquad
 T^*=Q^{2A}T,\qquad \Sigma_e=Q^{2A}\sigma_e,\\
 J_e^*=Q^{2A-(e+1)/2}J_e
             =\int_0^\infty R^e S^*_{z\tau_e}\,dR,\qquad
 M_e^*=Q^{2A-e/2}M_e=\varepsilon\partial_ZJ_e^*,\\
 q_*=q/Q,\qquad
 b_e^*=Q^{(e+1)/2}b_e(\sqrt Q R,z,t)
       =q_*^{-(e+1)/2}\widehat b_e(R/\sqrt{q_*}).
\end{gathered}
\end{equation}
The two moment powers follow directly from
$r^e\,dr=Q^{(e+1)/2}R^e\,dR$ and
$\partial_z=Q^{-D}\partial_Z$; the remaining factor in $M_e^*$
is $Q^{1/2-D}=\varepsilon$. In particular the scaled bump satisfies
$\int R^e b_e^*\,dR=1$.

Every term of the corrected primitive has the same exact stress unit:
\begin{equation}\label{mcstress:chart_sigma}
\begin{gathered}
 \Sigma_e=-R^{-e}\int_0^R X^e
             \bigl(F_e^*-\varepsilon\partial_Z(b_e^*J_e^*)\bigr)\,dX,\\
 (\partial_R+e/R)\Sigma_e
       =-F_e^*+\varepsilon\partial_Z(b_e^*J_e^*),\\
 \Sigma_e=\Sigma_e^{\rm old}
       -\frac{\varepsilon R\partial_Zq_*}{2q_*}b_e^*J_e^*.
\end{gathered}
\end{equation}
To verify the first line, insert
$F_e=Q^{-2A-1/2}F_e^*$ and $b_eJ_e=Q^{-2A}b_e^*J_e^*$ in
\eqref{mcstress:sigma}, and change the integration variable to
$s=\sqrt Q X$. The second follows by differentiation. For the third,
the physical correction in \eqref{mcstress:primitive_correction}
has factor $\sqrt Q\,q_z/q$; using
$q_z=Q^{1-D}\partial_Zq_*$ gives
$\sqrt Q\,q_z/q=\varepsilon\partial_Zq_*/q_*$.
This verifies the original $q$ derivative, including its sign.
Also, with $C_e^*=Q^{2A-(e+1)/2}C_e$,
\[
 \Sigma_e=-S^*_{r\tau_e}
                     -\varepsilon R^{-e}\partial_Z C_e^*.
\]

The completed stress and its pressure become exactly
\begin{equation}\label{mcstress:chart_tensor}
\begin{gathered}
 T^*_{rr}=0,\quad T^*_{r\theta}=\Sigma_2,\quad
 T^*_{rz}=\Sigma_1,\quad T^*_{z\theta}=-b_2^*J_2^*,\quad
 T^*_{zz}=-b_1^*J_1^*,\\
 T^*_{\theta\theta}=RF_r^*+\varepsilon R\partial_Z\Sigma_1,
 \qquad
 p_T^*=Q^{2A}p_T
  =\tfrac13(RF_r^*+\varepsilon R\partial_Z\Sigma_1-b_1^*J_1^*),\\
 (T^0)^*=T^*-p_T^*I_3,\qquad
 F^*+\operatorname{div}_*(T^0)^*+\nabla_*p_T^*=0.
\end{gathered}
\end{equation}
For example the radial row is
$-T^*_{\theta\theta}/R+\varepsilon\partial_Z\Sigma_1=-F_r^*$,
and the tangential rows use \eqref{mcstress:chart_sigma}. The pressure
identity follows from the same Cartesian isotropic-tensor calculation
as before. Physical divergence and pressure gradient both have factor
$Q^{-2A-1/2}$ because $\nabla_x=Q^{-1/2}\nabla_*$ on the chart.
Thus the radial residual, hoop stress, both surviving axial fluxes,
varying-bump correction and compact pressure all have their exact
source units. No decay or uniform estimate at a singular endpoint is
inferred from these finite compact-support identities.
